
\appendix
\renewcommand{\chaptername}{APPENDIX}
\titleformat{\chapter}[display]{\centering\bfseries\fontsize{16}{20}\selectfont}{APPENDIX \thechapter}{8pt}{\fontsize{16}{20}\selectfont\MakeUppercase}
\renewcommand{\thechapter}{\Alph{chapter}}
\renewcommand{\thefigure}{\thechapter.\arabic{figure}}
\renewcommand{\thetable}{\thechapter.\arabic{table}}

\clearpage
\chapter{Engine Constants and Thresholds}
\label{app:constants}
\begin{table}[!htbp]
\centering
\caption{Rule-layer thresholds (defaults, adjustable on the Incidents page).}
\small
\renewcommand{\arraystretch}{1.15}
\begin{tabularx}{\textwidth}{>{\raggedright\arraybackslash}p{5.4cm}>{\raggedright\arraybackslash}p{1.9cm}>{\raggedright\arraybackslash}X}
\toprule
\textbf{Name} & \textbf{Default} & \textbf{Meaning}\\
\midrule
\texttt{minHistorySamples} & 12 & Samples a service needs before it is judged\\
\texttt{zScoreAnomaly} & 3.5 & Deviation counted as anomalous (noise-normalised)\\
\texttt{zScoreHigh} & 5.0 & Deviation at or above this is HIGH\\
\texttt{zScoreCritical} & 8.0 & Deviation at or above this is CRITICAL\\
\texttt{minDeltaPercent} & 5 & Smallest change worth a flag\\
\texttt{absoluteHighPercent} & 70 & Absolute CPU or memory share counted as high\\
\texttt{absoluteCriticalPercent} & 85 & Absolute share counted as critical\\
\texttt{flatlineStdDev}, \texttt{flatline\-DeltaPercent} & 0.001, 30 & A series that goes flat after moving\\
{\ttfamily criticalServices\-ForCritical} & 2 & Critical services before the incident is CRITICAL\\
\texttt{silenceSamples} & 3 & Missed readings before a service counts as silent\\
\texttt{monitoringGapShare} & 0.8 & Share of silent services that means the monitoring itself is down\\
\texttt{explainAwayGamma} & 1.0 & Weight of what a sick callee explains ($\gamma$)\\
\texttt{cascadeEdgeProbability} & 0.5 & Starting propagation probability of an edge\\
\bottomrule
\end{tabularx}
\end{table}

\begin{table}[!htbp]
\centering
\caption{Edge evidence confidence.}
\small
\begin{tabular}{ll}
\toprule
\textbf{Evidence} & \textbf{Confidence $c(e)$}\\
\midrule
Seen in request logs, or oriented by a listening port & 1.00\\
Oriented only by the ephemeral-port range & 0.85\\
Observed, no direction evidence & 0.70\\
Observed, direction unresolved & 0.50\\
Declared in Compose, never seen carrying traffic & 0.40\\
\bottomrule
\end{tabular}
\end{table}

\clearpage
\chapter{Fault Classes}
\label{app:faults}
\begin{table}[!htbp]
\centering
\caption{Fault classes, how they are applied, and how many faults of each class the evaluation sets contain.}
\small
\renewcommand{\arraystretch}{1.2}
\begin{tabularx}{\textwidth}{>{\raggedright\arraybackslash}p{2.3cm}>{\raggedright\arraybackslash}X>{\centering\arraybackslash}p{1.6cm}>{\centering\arraybackslash}p{1.6cm}}
\toprule
\textbf{Class} & \textbf{Applied as} & \textbf{Real (62)} & \textbf{Synthetic (178)}\\
\midrule
crash & Process killed (\texttt{SIGKILL}); the container exits with a non-zero code & 16 & 27\\
stop & Container stopped with \texttt{docker stop} (exit code 0) & 16 & 24\\
brownout & Resource limit that starves the service without stopping it & 14 & 26\\
memory & Memory pressure on the service (a leak, or a limit) & 6 & 18\\
crash-loop & Repeated outages within minutes & 4 & 16\\
pause & Container paused & 4 & 15\\
flap & Repeated pause and release cycles & 2 & 12\\
latency & Slow responses & -- & 25\\
partition & Network detached from the service & -- & 15\\
\bottomrule
\end{tabularx}
\end{table}

\clearpage
\chapter{User Guide}
\label{app:guide}
\section*{Install and start}
\begin{lstlisting}[language=bash,numbers=none]
git clone https://github.com/Royson-salis-18/microservice-mapper.git
cd microservice-mapper
npm start        # installs what is missing, starts the server (:3001) and the dashboard (:5173)
\end{lstlisting}
Open \texttt{http://localhost:5173}. Node.js 20.19 or 22.12 and later is required. \texttt{npm run doctor} checks the machine.

\section*{Add a host}
Click ``Add project'' and fill in the target identifier, display name, the host's public IPv4 address, the SSH user and the path to the key file. The key stays on your machine. \texttt{npm run hosts} lists the configured hosts and tests port 22; \texttt{npm run hosts -- set <target> <ip>} changes an address.

\section*{Read an incident}
Open \emph{Incidents and cascade}. The headline names the most likely root cause with its confidence. The ranking shows what each service's score is made of; ``Why'' shows the own, symptom and walk parts and the evidence; the cascade panel shows who is already affected and who is next; the next steps are commands to run on the host.

\section*{Compare the three engines}
Open \emph{Detection engines}. The three columns show, for the same readings, what each engine detected, its ranked suspects and what it predicts. The switch in the top bar (or the ``Use'' buttons) chooses which engine decides the incident.

\section*{Generate load and reproduce results}
\begin{lstlisting}[language=bash,numbers=none]
npm test                 # the 339 tests
npm run eval             # record and score a campaign against the live mapper
npm run ml:compare       # detectors and root-cause methods on the recorded runs
npm run ml:hybrid        # fusion variants
npm run ml:predict       # early warning and cascade prediction
report/build.sh          # rebuild this report (needs pdflatex; Mermaid optional)
\end{lstlisting}

\section*{When nothing shows up}
\begin{lstlisting}[language=bash,numbers=none]
npm run check-host -- <public-ip> ~/path/to/key.pem 8080
npm run mapper-edges -- <project> <service-name-part>
\end{lstlisting}
The first says whether the host or the mapper is at fault; the second prints what the running mapper holds for a service.
